% === CODIFICACIÓN E IDIOMA ===
\usepackage[T1]{fontenc} % <--- AGREGADO: Mejora la guionización y corte de palabras en español
\usepackage[utf8]{inputenc}
\usepackage[spanish,es-tabla]{babel} % es-tabla cambia "Cuadro" por "Tabla"

% === MÁRGENES Y GEOMETRÍA ===
\usepackage[margin=2.5cm]{geometry}

% === ALINEACIÓN Y TEXTO ===
\usepackage{ragged2e} % <--- AGREGADO: Paquete necesario para el control de justificado profesional

% === GRÁFICOS Y FIGURAS ===
\usepackage{graphicx}
\graphicspath{{imagenes/}} % Define la carpeta raíz para las imágenes
\usepackage{booktabs} % Para tablas con estilo profesional
\usepackage{subcaption} % Para poner múltiples imágenes juntas

% === MATEMÁTICAS AVANZADAS ===
\usepackage{amsmath,amsfonts,amssymb}

% === HERRAMIENTAS ADICIONALES ===
\usepackage{url} % Manejo correcto de direcciones web
\usepackage{cite} % Agrupa citas consecutivas de forma limpia (ej: [1-3])
\usepackage{hyperref} % Añade hipervínculos interactivos en el PDF (Índice, citas)
\hypersetup{
    colorlinks=true,
    linkcolor=blue,
    filecolor=magenta,      
    urlcolor=cyan,
    citecolor=green
}

% === NOTAS DE SEGUIMIENTO (Opcional para el equipo) ===
\usepackage{todonotes} % Permite usar \todo{Falta revisar esto} en el texto

\usepackage{titlesec}

% Modifica el espaciado de \chapter
% [izquierda] {espacio_superior} {espacio_entre_titulo_y_texto}
\titlespacing*{\chapter}{0pt}{-20pt}{20pt}